% generated by atam/source/part_scorecard.py - do not edit by hand
\paragraph{Risks}
\begin{itemize}
\item \textbf{R-6 The comparison engine's daily calibration amplifies small data changes.} On the faulty feed 11.9\% of Realized-CAViaR's single-stock forecasts move more than 5\% (without the jump filter: 15.2\%). On 12\% of the stock-days the daily calibration of the two feeds lands on different optima (a beta more than 0.05 apart), and 54\% of the forecasts off fall on those days: |R| and RV of the same hour are close substitutes, so the pinball loss is nearly flat along a direction of (beta2, beta3) and a few changed observations move the optimum, and with it the whole day's VaR. The calibration starts from the previous day's parameters, which keeps it when it still fits best; a rule that keeps it unless the new window rejects it froze the parameters and made the two feeds diverge more (tested and rejected). The faults reach the model through its estimates, not its inputs: the protection that works for the portfolio engines is not enough for a model fitted on 40 sessions of quantile losses.
\item \textbf{R-7 The single-instrument engine is timed at a different boundary.} EWMA and GARCH-t time the same bar cycle of a 15-stock book (fills, the shared jump filter, every model, VaR; the option revaluation runs off the bar thread and is reported apart). Realized-CAViaR times the model stage of one instrument in batches of 1000, whose percentiles are of batch means, so its numbers are not comparable with theirs: the cost side of the research question must use the matched model-stage benchmark, where the three models share one boundary.
\item \textbf{R-9 Moderate doses of every recoverable fault still move the VaR.} On faulty copies of both datasets with missing, duplicated, out-of-order and late rows, halts, frozen quotes, bad ticks, corrupt values and early-ending sessions together, Realized-CAViaR stays above 10\% of forecasts more than 5\% away from the clean run (EWMA: synthetic 0.7\%, b3\_real 6.3\% of the portfolio forecasts more than 5\% off, 0.3\%, 1.4\% not issued; GARCH-t: synthetic 2.5\%, b3\_real 9.5\% of the portfolio forecasts more than 5\% off, 0.3\%, 1.4\% not issued; Realized-CAViaR: synthetic 22.5\%, b3\_real 22.5\% of the single-stock forecasts (every stock) more than 5\% off, 2.4\%, 3.6\% not issued). atam/results/faults/ shows which fault does it, one kind at a time: missing data is the residual problem (no protection can recover a print that never arrives).
\end{itemize}
\paragraph{Sensitivity points}
\begin{itemize}
\item \textbf{S-3 The matched comparison does not cover the opening hour.} CAViaR needs one hour of returns of the same session, so Realized-CAViaR scores every model on targets starting 11:00-16:45 only. The most volatile bars of the day (the opening hour) are outside the comparison, so its conclusions about the models hold from 11:00 on.
\item \textbf{S-4 Reuse by compiling the other modules' sources.} Besides the shared library common/ (model interface, baseline, jump filter, with its own tests), Realized-CAViaR compiles 10 source files of ewma and garch\_t (loader, EWMA, GARCH, backtest, the jump filter's panel adapter). It guarantees the same conventions and estimators in the comparison, but couples the modules with no version: a change in ewma or garch\_t silently changes the comparison engine. atam/run.bat rebuilds it before every evaluation, so a rerun after such a change is the safeguard.
\item \textbf{S-6 An unprinted bar widens the next forecast on purpose.} On the faulty feed 1.46\% of the GARCH-t forecasts move more than 5\% from the clean-feed ones, but 38 of them are issued at or right after a bar in which a stock did not print; without those the share is 0.23\%. The engine marks a stock that skips a bar at its last print, so the P\&L of its next bar spans two bars, and it adds the variance of the unprinted bar to that forecast: the clean-feed VaR covers one bar, the faulty-feed one covers the P\&L it is actually scored on. R1 counts this deliberate response as a deviation; the threshold would be met if the scenario compared forecasts of the same P\&L horizon.
\end{itemize}
\paragraph{Trade-off points}
\begin{itemize}
\item \textbf{T-1 Data protection: robustness vs reaction to genuine shocks.} With the EWMA data-quality layer and the jump filter, 0.00\% of the faulty-feed forecasts stray more than 5\% from the clean-feed ones (without them: 10.01\%). The price is paid on genuine moves: at a shock bar the VaR moves up to 3.9\% away from the unprotected one, for 27 bar(s), because a move confined to one stock is held until the next bar confirms it. Timed latency cost: +0.60 us at the median per bar (the filter runs inside the timed bar).
\item \textbf{T-2 GARCH-t: a stateful jump filter against stateless clipping.} On the faulty feed, 25.6\% of the unprotected GARCH-t forecasts stray more than 5\%; clipping shocks at 6 sd leaves 16.0\% and the jump filter (default) 1.46\%. With every recoverable fault together (R3): 66.7\% unprotected, 38.0\% clipped, 9.5\% filtered. At the genuine shocks of the clean feed, the lowest ratio to the unprotected VaR is 0.973 with clipping and 0.997 with the filter: the filter delays a move confined to one stock by one bar, then lets it in at full size, while clipping shrinks every large shock for as long as it is in the variance.
\end{itemize}
\paragraph{Non-risks}
\begin{itemize}
\item \textbf{N1 No look-ahead in the data path.} Replacing every price after a cut time leaves all earlier forecasts and calibrations unchanged in EWMA, GARCH-t and Realized-CAViaR (1,250,180 cells over three cuts); the same test changes 100\% of the later forecasts (lowest engine), so it would have caught a leak.
\item \textbf{N2 Replays are deterministic.} Three runs of EWMA, GARCH-t and Realized-CAViaR produce identical output files apart from timing columns (1,015,686 cells compared), including the multithreaded GARCH-t refits, the Monte Carlo option revaluation and the multi-start CAViaR fits.
\item \textbf{N3 The portfolio engines read the market identically on clean data.} With both data protections off, the EWMA(0.94) baseline that EWMA and GARCH-t both compute is identical forecast by forecast on the synthetic market: the two pipelines read, align and return the market in the same way.
\item \textbf{N5 The model comparison runs on the shared market and on matched targets.} Realized-CAViaR reads the market with the repository's shared convention (synthetic: 2638 of 2638 target returns equal the shared convention; b3\_real: 1120 of 1120 target returns equal the shared convention; the three models share one target record per forecast by construction). It is the engine in which the research question can be answered fairly.
\item \textbf{N4 Latency and recalibration are far below the bar interval.} Per-bar updates of the portfolio engines take microseconds (EWMA p99 2.9 us, GARCH-t p99 5.0 us; Realized-CAViaR updates in 16 ns at the model stage) and the daily refit at most EWMA 6 ms, GARCH-t 106 ms, Realized-CAViaR 38 ms, against 15-minute bars: at this scale performance does not constrain the designs. Timings vary between sessions on the same machine (up to a factor of two), so verdicts near a threshold are marked PARTIAL.
\item \textbf{N6 The main model is the cheapest one under a common boundary.} In the matched model-stage benchmark (Realized-CAViaR 16 ns / 40 B, GARCH-t 23 ns / 280 B, EWMA-t 25 ns / 248 B, EWMA-N 25 ns / 240 B) Realized-CAViaR has the lowest update cost and the smallest state: the computational side of the research question does not penalize it.
\item \textbf{N8 One data protection and one baseline in every engine.} The jump filter (common/, causal, decided on each bar with the bars before it) runs by default in EWMA, GARCH-t and Realized-CAViaR, and the RiskMetrics baseline is one implementation behind the common model interface in both portfolio engines: the baseline coincides forecast by forecast on the synthetic market (0 of the forecasts differ; with both protections off 0), and on the B3 data 0 differ by default, 0 with both off.
\item \textbf{N9 One model interface for the portfolio engines.} In EWMA and GARCH-t every VaR model implements riskcommon::PortfolioModel and the replay, the backtest and the output files iterate over the registered models: a new model is a class in the model layer plus its registration (EWMA: 2 file(s) outside the layer name a model, GARCH-t: 1 file(s) outside the layer name a model). Realized-CAViaR keeps its models in one orchestrator, selected by position (2 files to change), outside the class interface: its models are univariate quantile filters with another life cycle.
\item \textbf{N10 The target time grid is a configuration in every engine.} Realized-CAViaR has a high-frequency mode (5-second prices, five 1-minute phases of a 5-minute horizon); the portfolio engines aggregate 5-second samples into the bars of any --session grid (--resample), so a 5-minute VaR every minute is five runs on 5-minute grids opening one minute apart, each with non-overlapping targets as CAViaR's phases. The experiments of the monograph run on the shared 15-minute data, which is the grid they report.
\item \textbf{N11 The option revaluation stays off the bar update.} The Monte Carlo revaluation of the option book (the revaluation itself takes 6.9 ms per bar at the median) runs after the VaR is published, off the bar thread, so the bar update keeps its own cost: GARCH-t p99 5.0 us with the book, 5.0 us without it (p50 2.3 vs 2.2 us); its records are the same as with the revaluation on the bar thread (unit test).
\end{itemize}
